where @@M0@@ is the structural morphism @@M1@@.

The @@M2@@ considered will be those defined using locally free sheaves of finite type. We shall denote by @@M3@@, @@M4@@, @@M5@@ the classes of @@M6@@, @@M7@@, @@M8@@. Put @@M9@@, @@M10@@, and denote by @@M11@@, @@M12@@ the classes of @@M13@@, @@M14@@ in @@M15@@. Since @@M16@@ and @@M17@@ belong to @@M18@@, we have @@M19@@, i.e. @@M20@@ mod. @@M21@@ i.e. @@M22@@

\sgaStableTarget{sga6.e06.sec.5.2}5.2.
If @@M23@@ is an @@M24@@-Module locally free of finite type, the @@M25@@ are in general not locally free of finite type on @@M26@@, so that a homomorphism @@M27@@ cannot be defined as it is for the @@M28@@ defined using perfect complexes. Theorem 1.1 allows us to get around this difficulty as follows: if @@M29@@, there is a unique way to write:
@@M30@@
we then put:
@@M31@@

Clearly the homomorphism thus defined is additive and satisfies the projection formula, i.e. is @@M32@@-linear.

The following relations hold:
@@M33@@

To verify the last group of relations, we proceed by induction on @@M34@@. Suppose that, for @@M35@@, @@M36@@; relation (1.9.1), multiplied by a suitable power of @@M37@@, gives:
@@M38@@
We can apply @@M39@@ to it:
